\section{Introduction}
\label{sec:introduction}

% --- Motivation ---------------------------------------------------------
LLM-based assistants have become part of everyday software development, generating
anything from small snippets to complete functions. Benchmarks for code generation
predominantly measure \emph{functional correctness}, e.g., the share of generated
solutions that pass a set of unit tests. Non-functional properties such as runtime
performance and, in particular, energy efficiency are rarely considered, even though
generated code may be executed many times in production and on a wide range of
devices, from data-centre servers to resource-constrained embedded systems.
\TODO{Add 1--2 citations on the ICT/software share of energy consumption and on the
adoption of AI coding assistants.}

% --- Problem ------------------------------------------------------------
Early studies indicate that the energy efficiency of LLM-generated code varies
between models and often falls short of human-written reference
solutions~\cite{solovyeva2025powerhungry,islam2025evaluating,apsan2026whereare}.
However, existing studies typically vary only a subset of the factors that plausibly
influence energy consumption: they either cover several languages on a small number of
similar platforms~\cite{solovyeva2025powerhungry}, or several heterogeneous platforms
for a single language~\cite{apsan2026whereare}. It therefore remains unclear whether a
statement such as ``model~A generates more energy-efficient code than model~B'' holds
in general, or only for a particular language and execution environment.

% --- Goal & approach ----------------------------------------------------
This paper addresses this question with a full factorial experiment in which
three LLMs, accessed through their vendors' official coding CLIs, generate solutions for
14~programming tasks in Python, Java, and C++.
All functionally correct solutions are executed on three heterogeneous platforms
(an x86 laptop, an x86 server, and an ARM-based Raspberry~Pi), and their energy
consumption, execution time, and memory usage are measured under controlled conditions.
Our central research question is:

\begin{quote}
\emph{To what extent do the choice of LLM, programming language, and hardware platform
influence the energy efficiency of LLM-generated code?}
\end{quote}

% --- Contributions ------------------------------------------------------
The contributions of this paper are:
\begin{itemize}[leftmargin=*]
  \item A full factorial empirical study of LLM-generated code across
        three LLMs (Claude, GPT, Gemini via their official CLIs), three programming
        languages, and three hardware platforms.
  \item A statistical analysis of main \emph{and interaction} effects, with a particular
        focus on whether the energy-efficiency ranking of LLMs is stable across
        languages and platforms.
  \item An analysis of the relationship between energy, execution time, and memory for
        generated code.
  \item A replication package containing tasks, prompts, generated code, measurement
        scripts, and raw data.
\end{itemize}

\TODO{Paper outline paragraph (``The remainder of this paper is structured as
follows\ldots'') once the structure is final.}
